\documentclass[11pt]{article}

% ===================== PACKAGES =====================
\usepackage[a4paper,margin=1in]{geometry}
\usepackage{amsmath,amssymb}
\usepackage{enumitem}
\usepackage{fancyhdr}
\usepackage{xcolor} 
\usepackage{algorithm}
\usepackage{algpseudocode}

% ===================== HEADER & FOOTER =====================
\pagestyle{fancy}
\fancyhf{}
\lhead{CSE 400: Fundamentals of Probability in Computing}
\rhead{Lecture 10 Scribe}
\cfoot{\thepage}

% ===================== TITLE =====================
\title{
    \normalsize School of Engineering and Applied Science (SEAS), Ahmedabad University \\
    \vspace{0.2cm}
    \textbf{CSE 400: Fundamentals of Probability in Computing}\\
    \Large Lecture 10 Scribe: Randomized Min-Cut Algorithm
}
\author{Instructor: Dhaval Patel, PhD} 
\date{}

\begin{document}
\maketitle

\vspace{-1cm}
\begin{center}
    \begin{tabular}{ll}
        \textbf{Date:} & {February 5, 2026 \hspace{2.2in}} \\ [1.5ex]
        
    \end{tabular}
\end{center}

\hrule
\vspace{0.5cm}

% ===================== CONTENT =====================

\section{Outline of the Lecture}
The lecture covers the following topics:
\begin{itemize}
    \item Min-Cut Problem (Definition and Applications)
    \item Max-Flow Min-Cut Theorem
    \item Deterministic Min-Cut Algorithm (Stoer--Wagner)
    \item Randomized Min-Cut Algorithm (Karger’s Algorithm)
    \item Comparison: Deterministic vs Randomized Approaches
    \item Python Simulation
\end{itemize}

\section{Min-Cut Problem}

\subsection{Why Use Min-Cut?}
We use the min-cut algorithm in various applications to solve problems related to network connectivity, reliability, and optimization.

\subsubsection*{1. Network Design}
Min-cut helps in improving the efficiency of communication and optimizing network flow.
\begin{itemize}
    \item \textbf{Objective:} Improve communication efficiency.
    \item \textbf{Objective:} Optimize network flow.
    \item \textbf{Method:} Identify the minimum capacity cut.
\end{itemize}

\subsubsection*{2. Communication Networks}
Minimum cut is useful for understanding the vulnerability of networks to failures.
\begin{itemize}
    \item Evaluates network vulnerability.
    \item Assists in building robust and fault-tolerant communication networks.
\end{itemize}

\subsubsection*{3. VLSI Design}
In Very Large Scale Integration (VLSI) design, the algorithm is used for partitioning circuits into smaller components.
\begin{itemize}
    \item Leads to reduced interconnectivity complexity between components.
\end{itemize}

\section{What is Min-Cut?}

\subsection{Definition: Cut-Set}
A \textbf{cut-set} in a graph is a set of edges whose removal breaks the graph into two or more connected components.

\subsection{Definition: Minimum Cut}
Given a graph $G = (V, E)$ where $V$ is the set of vertices and $E$ is the set of edges. The \textbf{minimum cut (min-cut)} problem is to find a minimum cardinality cut-set in $G$.
\begin{itemize}
    \item Among all possible cut-sets, select the one with the smallest number of edges.
\end{itemize}

\subsection{Sensitivity of Randomized Algorithms}
Min-cut algorithms (like Karger’s) are random and sensitive to the initial choice of edges. If the algorithm contracts critical edges early in the process, it might find a smaller cut (or fail to find the true min-cut).

\section{Edge Contraction}
The main operation in the randomized algorithm is \textbf{edge contraction}.
Let the edge be $(u, v)$. Contraction consists of:
\begin{enumerate}
    \item Merging vertices $u$ and $v$ into one vertex.
    \item Eliminating all edges connecting $u$ and $v$ (self-loops).
    \item Retaining all other edges (resulting in potential parallel edges).
\end{enumerate}

\section{Successful vs. Unsuccessful Runs}
\begin{itemize}
    \item \textbf{Successful Min-Cut Run:} An execution of the algorithm where it correctly identifies the minimum cut of the graph.
    \item \textbf{Unsuccessful Min-Cut Run:} An iteration where the algorithm fails to correctly identify the minimum cut (e.g., due to contracting an edge that belongs to the min-cut set).
\end{itemize}

\section{Max-Flow Min-Cut Theorem}
\textbf{Theorem Statement:} In a flow network, the maximum amount of flow passing from the source to the sink is equal to the total weight of the edges in a minimum cut.
\[ \text{Maximum Flow} = \text{Minimum Cut Capacity} \]

\subsection*{Definitions}
\begin{itemize}
    \item \textbf{Capacity of a cut:} The sum of the capacities of the edges in the cut oriented from set $X$ to set $Y$.
    \item \textbf{Minimum cut:} The cut in the network with the smallest possible capacity.
    \item \textbf{Maximum flow:} The largest possible flow from source $S$ to sink $T$.
\end{itemize}

\section{Deterministic Min-Cut Algorithm}
\subsection{Stoer--Wagner Minimum Cut Algorithm}
Let $s$ and $t$ be two vertices of graph $G$. Let $G/\{s, t\}$ be the graph obtained by merging $s$ and $t$.
A minimum cut of $G$ can be obtained by taking the smaller of:
\begin{enumerate}
    \item A minimum $s-t$ cut of $G$.
    \item A minimum cut of $G/\{s, t\}$.
\end{enumerate}

\subsection{Pseudocode (Deterministic)}
\begin{algorithm}
\caption{MinimumCutPhase(G, a)}
\begin{algorithmic}
\State $A \leftarrow \{a\}$
\While{$A \neq V$}
    \State Add to $A$ the most tightly connected vertex
\EndWhile
\State \Return The cut weight as the ``cut of the phase''
\end{algorithmic}
\end{algorithm}

\begin{algorithm}
\caption{MinimumCut(G)}
\begin{algorithmic}
\While{$|V| \ge 1$}
    \State Choose any $a$ from $V$
    \State MinimumCutPhase(G, a)
    \If{cut-of-the-phase $<$ current minimum cut}
        \State Store cut-of-the-phase as current minimum cut
    \EndIf
    \State Shrink $G$ by merging the two vertices added last
\EndWhile
\State \Return The minimum cut
\end{algorithmic}
\end{algorithm}

\section{Randomized Min-Cut Algorithm}
\subsection{Why Randomized Algorithm?}
Randomized algorithms provide:
\begin{itemize}
    \item A probabilistic guarantee of success.
    \item A more accurate estimate of the minimum cut with fewer iterations.
    \item Efficiency and ease of parallelization.
\end{itemize}

\subsection{Karger's Algorithm Theorem}
The algorithm outputs a min-cut set with probability at least:
\[ \frac{2}{n(n-1)} \]
where $n$ is the number of vertices.

\section{Comparison: Deterministic vs Randomized}

\begin{table}[h]
\centering
\begin{tabular}{|l|p{6cm}|p{6cm}|}
\hline
\textbf{Feature} & \textbf{Deterministic (Stoer-Wagner)} & \textbf{Randomized (Karger)} \\
\hline
\textbf{Guarantee} & Always guarantees an exact minimum cut. & Provides an approximate minimum cut with high probability. \\
\hline
\textbf{Complexity} & $O(V \cdot E + V^2 \log V)$ & $O(V^2)$ (per run) \\
\hline
\end{tabular}
\end{table}

\section{Python Simulation}
\begin{itemize}
    \item Open the Campuswire Post regarding Lecture 10.
    \item Download the \texttt{.ipynb} file pasted in the comments to run the simulation.
\end{itemize}

\bigskip
\hrule
\vspace{0.2cm}
\begin{center}
    \small \textit{End of Lecture 10 Scribe}
\end{center}

\end{document}